\chapter{Założenia i architektura aplikacji}
\label{chap:projekt-programu}

\section{Założenia projektowe}
\label{sec:zalozenia-projektowe}

Celem części programistycznej pracy jest przygotowanie aplikacji, która pozwala przejść od prostego opisu tekstowego obwodu do wyniku analizy numerycznej. Program nie jest traktowany jako uniwersalny zamiennik rozbudowanych środowisk SPICE, lecz jako dedykowane narzędzie pokazujące, w jaki sposób metoda ZMPW może zostać zaimplementowana w natywnej aplikacji napisanej w języku Swift.

Najważniejsze założenie projektowe polega na oddzieleniu opisu obwodu od sposobu jego rozwiązania. Użytkownik może wpisać układ bezpośrednio w głównym oknie programu albo przeciągnąć do niego plik tekstowy z netlistą. Aplikacja przekształca ten opis na obiekty reprezentujące stemple elementów. Dopiero tak przygotowana lista stempli jest przekazywana do solvera, który buduje macierz globalną i wyznacza wynik wybranej analizy.

Takie podejście pozwala rozwijać program stopniowo. Dodanie nowego elementu wymaga zdefiniowania jego stempla oraz rozszerzenia warstwy parsera o regułę rozpoznawania nowego oznaczenia. Sam solver nie musi znać nazw elementów ani formatu pliku wejściowego, ponieważ operuje wyłącznie na macierzach zwracanych przez stemple.

\section{Ogólna architektura}
\label{sec:ogolna-architektura}

Program został podzielony na kilka warstw odpowiedzialności. Pierwszą z nich jest interfejs użytkownika, który udostępnia edytor opisu obwodu, obsługę przeciągania plików \texttt{.txt} oraz przycisk uruchamiający symulację. Drugą warstwą jest parser tekstu wejściowego. Jego zadaniem jest odczytanie kolejnych linii, rozpoznanie typu elementu, sprawdzenie liczby parametrów oraz utworzenie listy obiektów typu \texttt{Stamp}. Trzecią warstwę tworzą klasy elementów obwodu, na przykład rezystor, źródła niezależne i źródła sterowane. Każdy z tych elementów zwraca lokalny wkład do macierzy przewodności oraz do wektora wymuszeń. Czwartą warstwą jest solver, który sumuje wszystkie stemple i rozwiązuje powstały układ równań.

Na poziomie przepływu danych aplikacja wykonuje następujące kroki:

\begin{enumerate}
  \item pobiera tekstowy opis obwodu z edytora albo z przeciągniętego pliku,
  \item dzieli opis na linie i tokeny,
  \item rozpoznaje typy elementów oraz komendę analizy,
  \item tworzy obiekty stempli i przypisuje dodatkowe indeksy równań,
  \item przekazuje listę stempli oraz komendę do solvera,
  \item buduje macierz globalną i oblicza wynik analizy,
  \item prezentuje wynik punktu pracy albo wykres przebiegów czasowych.
\end{enumerate}

Wynikiem pracy parsera jest struktura zawierająca tablicę stempli, komendę solvera oraz listę węzłów, których przebiegi mają zostać pokazane na wykresie. Dzięki temu parser nie rozpoczyna obliczeń samodzielnie. Jego odpowiedzialność kończy się na przygotowaniu danych wejściowych w formie zrozumiałej dla dalszych części programu.

Cały przepływ danych można przedstawić jako przejście od interfejsu użytkownika albo pliku tekstowego do struktury \texttt{ParserResult}, a następnie do solvera i modułu wizualizacji. Dla komendy \texttt{.op} solver zwraca pojedynczy wektor rozwiązania. Dla komendy \texttt{.tran} zwraca sekwencję próbek czasowych, które moduł \texttt{Plotter} zamienia na serie punktów dla węzłów wskazanych przez \texttt{.show}. Interfejs użytkownika nie interpretuje samodzielnie składni netlisty, lecz przekazuje tekst do parsera i wybiera odpowiedni widok wyniku na podstawie zwróconej komendy.

\Cref{fig:main-view} pokazuje główne okno aplikacji. Centralną część zajmuje edytor netlisty, a po prawej stronie znajduje się obszar do przeciągania plików tekstowych. Po uruchomieniu symulacji program wybiera sposób prezentacji wyniku na podstawie komendy analizy: dla \texttt{.op} pokazuje wartości niewiadomych, a dla \texttt{.tran} rysuje wykres napięć wskazanych komendą \texttt{.show}.

\begin{figure}[H]
\centering
\includegraphics[width=0.9\textwidth]{MainView.png}
\caption{Główne okno aplikacji z edytorem opisu układu}
\label{fig:main-view}
\end{figure}

\section{Format pliku wejściowego}
\label{sec:format-wejsciowy}

Opis obwodu ma postać tekstową zbliżoną do uproszczonej netlisty SPICE. Jedna linia opisuje jeden element albo jedną komendę analizy. W przypadku elementu pierwszy token określa jego typ, kolejne tokeny oznaczają numery węzłów, a ostatnie wartości są parametrami elementu. Węzeł o numerze \texttt{0} pełni rolę węzła odniesienia. Taki format jest prosty do edycji ręcznej i dobrze pasuje do programowego parsowania opisu obwodu \cite{nagel1975spice2,kundert1995designer}.

Wartości liczbowe mogą być zapisane w postaci pełnej albo z przedrostkiem mnożącym. Parser rozpoznaje przedrostki z uwzględnieniem wielkości liter: \texttt{G} oznacza \(10^9\), \texttt{M} oznacza \(10^6\), \texttt{k} oznacza \(10^3\), \texttt{m} oznacza \(10^{-3}\), \texttt{u} oznacza \(10^{-6}\), \texttt{n} oznacza \(10^{-9}\), a \texttt{p} oznacza \(10^{-12}\). Dzięki temu zapis \texttt{1k}, \texttt{100n} albo \texttt{4m} jest równoważny odpowiednio wartościom \texttt{1000}, \texttt{0.0000001} i \texttt{0.004}. Rozróżnianie wielkości liter jest istotne, ponieważ \texttt{M} oznacza mega, natomiast \texttt{m} oznacza mili.

Aktualnie obsługiwane są: rezystor \texttt{R}, kondensator \texttt{C}, niezależne źródło prądowe \texttt{DCCS}, niezależne źródła napięciowe \texttt{DCVS} i \texttt{ACVS}, źródła sterowane \texttt{VCCS}, \texttt{VCVS}, \texttt{CCCS}, \texttt{CCVS} oraz tranzystor bipolarny \texttt{BJT}. Komenda \texttt{.op} oznacza analizę punktu pracy, natomiast \texttt{.tran} uruchamia analizę czasową. W przypadku analizy czasowej można dodatkowo użyć komendy \texttt{.show}, która wskazuje węzły przeznaczone do wizualizacji.

\begin{table}[H]
\centering
\caption{Format opisu obsługiwanych elementów}
\label{tab:format-elementow}
\begin{tabularx}{\textwidth}{@{}l>{\ttfamily\small}lX@{}}
\toprule
Oznaczenie & Postać linii & Znaczenie parametrów \\
\midrule
\texttt{R} & R n1 n2 R & rezystor między węzłami \texttt{n1} i \texttt{n2} \\
\texttt{C} & C n1 n2 C & kondensator między węzłami \texttt{n1} i \texttt{n2} \\
\texttt{DCCS} & DCCS n+ n- I & niezależne źródło prądowe DC \\
\texttt{DCVS} & DCVS n+ n- V & niezależne źródło napięciowe DC \\
\texttt{ACVS} & ACVS n+ n- A f & sinusoidalne źródło napięciowe o amplitudzie \texttt{A} i częstotliwości \texttt{f} \\
\texttt{VCCS} & VCCS in+ in- out+ out- gm & źródło prądowe sterowane napięciem \\
\texttt{VCVS} & VCVS in+ in- out+ out- A & źródło napięciowe sterowane napięciem \\
\texttt{CCCS} & CCCS in+ in- out+ out- beta Vx & źródło prądowe sterowane prądem pomocniczego źródła napięciowego \\
\texttt{CCVS} & CCVS in+ in- out+ out- Rm Vx & źródło napięciowe sterowane prądem pomocniczego źródła napięciowego \\
\texttt{BJT} & BJT b c e Is beta m & tranzystor bipolarny opisany węzłami bazy, kolektora i emitera \\
\bottomrule
\end{tabularx}
\end{table}

\Cref{lst:all-elements-netlist} przedstawia przykładową netlistę zawierającą wszystkie obsługiwane obecnie typy elementów. Sam układ nie musi mieć znaczenia praktycznego; jego zadaniem jest pokazanie składni wejściowej i kolejności parametrów.

\begin{lstlisting}[language=Netlist,caption={Przykładowy opis obwodu z wszystkimi obsługiwanymi elementami},label={lst:all-elements-netlist}]
R 1 0 1k
DCCS 0 1 1m
DCVS 2 0 5
VCCS 1 0 3 0 0.5
VCVS 1 0 4 0 2
CCCS 2 0 5 0 3 1
CCVS 3 0 6 0 4 1
C 1 2 100n
ACVS 7 0 2 1k
BJT 2 3 4 1p 100 1
.op
\end{lstlisting}

Dla źródeł napięciowych przyjęto konwencję zgodną z równaniami ZMPW: zapis \texttt{DCVS n+ n- V} oznacza warunek \(V_{n+} - V_{n-} = V\). Analogicznie \texttt{ACVS n+ n- A f} oznacza wymuszenie sinusoidalne \(V_{n+} - V_{n-} = A\sin(2\pi f t)\). Konwencja ta jest istotna, ponieważ od kolejności węzłów zależy znak wymuszenia.

\Cref{lst:emitter-follower-netlist} pokazuje netlistę układu wtórnika emiterowego użytego do sprawdzenia analizy czasowej. Węzeł \texttt{1} jest wyjściem źródła sinusoidalnego, węzeł \texttt{2} odpowiada bazie tranzystora, węzeł \texttt{3} kolektorowi, a węzeł \texttt{4} emiterowi. Komenda \texttt{.show 2 4} oznacza, że na wykresie mają zostać pokazane napięcia bazy i emitera.

\begin{lstlisting}[language=Netlist,caption={Opis układu wtórnika emiterowego},label={lst:emitter-follower-netlist}]
R 3 2 27k
R 0 4 1k
ACVS 1 0 2 1k
DCVS 3 0 5
C 1 2 100n
BJT 2 3 4 1p 100 1
.tran 4m 10u
.show 2 4
\end{lstlisting}

\section{Parser netlisty}
\label{sec:parser-netlisty}

Parser odczytuje plik linia po linii. Każda linia jest dzielona na tokeny rozdzielone białymi znakami. Linie puste oraz komentarze rozpoczynające się od znaku \texttt{*} są pomijane. Jeżeli pierwszy token jest komendą analizy, parser zapisuje ją jako polecenie dla solvera. W przeciwnym przypadku token jest interpretowany jako oznaczenie elementu.

Rozpoznawanie elementów jest wydzielone do typu \texttt{ElementType}. Typ ten pełni rolę słownika wspieranych elementów: zamienia tekstowe oznaczenie z pliku na konkretny wariant programu, określa wymaganą liczbę tokenów oraz wskazuje, które tokeny są numerami węzłów. Dzięki temu podstawowe reguły składniowe są skupione w jednym miejscu.

Parser wykonuje dwa przejścia po danych wejściowych. W pierwszym przejściu sprawdza strukturę linii, zapamiętuje poprawnie rozpoznane elementy i wyznacza największy numer węzła występujący w pliku. Drugie przejście służy do tworzenia stempli. Podział na dwa etapy jest potrzebny, ponieważ elementy zawierające dodatkowe równania, na przykład źródła napięciowe, wymagają indeksów większych od numerów zwykłych węzłów.

\begin{lstlisting}[language=Swift,caption={Uproszczony schemat odpowiedzialności parsera},label={lst:parser-flow}]
let parserResult = try Parser().parse(fileURL)
let stamps = parserResult.stamps
let command = parserResult.command
let result = try Solver().solve(stamps, command)
\end{lstlisting}

\section{Przypisywanie dodatkowych równań}
\label{sec:dodatkowe-rownania-program}

W metodzie ZMPW część elementów wymaga wprowadzenia dodatkowej niewiadomej. Dotyczy to przede wszystkim źródeł napięciowych oraz źródeł sterowanych prądem. W implementacji taka niewiadoma jest reprezentowana przez dodatkowy indeks macierzy, nazywany w programie \texttt{newRow}.

Parser najpierw zlicza węzły układu, a następnie nadaje kolejne indeksy elementom wymagającym dodatkowych równań. Jeżeli największy numer węzła w pliku wynosi \texttt{nodeCount}, pierwszy dodatkowy indeks ma wartość \texttt{nodeCount + 1}, drugi \texttt{nodeCount + 2} i tak dalej. Źródło \texttt{CCVS} potrzebuje dwóch takich indeksów, ponieważ oprócz równania wyjściowego źródła napięciowego wymaga również równania umożliwiającego reprezentację prądu sterującego.

Ten mechanizm jest istotny, ponieważ użytkownik nie powinien ręcznie numerować równań pomocniczych. Numery w pliku wejściowym opisują wyłącznie węzły obwodu, natomiast wewnętrzna rozbudowa macierzy jest zadaniem programu.

\section{Reprezentacja elementów i stempli}
\label{sec:reprezentacja-elementow}

Każdy element obwodu jest reprezentowany przez klasę albo strukturę zwracającą stempel. Stempel składa się z lokalnej macierzy przewodności oraz lokalnego wektora wymuszeń. Dla rezystora niezerowe wartości trafiają wyłącznie do macierzy przewodności. Dla niezależnego źródła prądowego wkład pojawia się w wektorze wymuszeń. Źródła napięciowe oraz część źródeł sterowanych wprowadzają dodatkowe wiersze i kolumny w macierzy globalnej.

Istotne jest to, że obiekty elementów nie znają całej netlisty. Otrzymują jedynie numery węzłów, parametry elementu oraz ewentualne indeksy dodatkowych równań. Następnie zwracają macierze, które można dodać do macierzy globalnej. Dzięki temu kod odpowiedzialny za konkretny element pozostaje lokalny i łatwy do testowania.

\section{Solver i komendy analizy}
\label{sec:solver-komendy}

Solver otrzymuje listę stempli oraz komendę analizy. Dla analizy punktu pracy \texttt{.op} tworzy macierz globalną, dodaje do niej wkład każdego stempla i rozwiązuje układ równań liniowych. Wynikiem jest macierz wartości niewiadomych, czyli napięć węzłowych oraz prądów związanych z dodatkowymi równaniami.

Komenda \texttt{.tran} wykorzystuje tę samą ideę składania stempli, ale obliczenia są wykonywane w kolejnych chwilach czasu. Elementy dynamiczne są wtedy zastępowane równoważnymi modelami dyskretnymi zależnymi od kroku czasowego i historii poprzednich stanów. Dzięki wspólnej reprezentacji przez stemple analiza czasowa może korzystać z tej samej warstwy macierzy co punkt pracy \cite{nagel1975spice2,vlach1994computer}. Wynik analizy czasowej może zostać przekazany do modułu wizualizacji, który tworzy serie punktów dla wskazanych węzłów.

\Cref{fig:op-result-view} pokazuje wynik analizy punktu pracy dla prostego dzielnika napięcia. Opis układu zawiera dwa rezystory \SI{1}{\kilo\ohm} oraz źródło napięciowe \SI{5}{\volt}. Ponieważ rezystory mają takie same wartości, napięcie w węźle pośrednim wynosi połowę napięcia zasilania, czyli około \SI{2.5}{\volt}. Wynik programu pokazuje również \SI{5}{\volt} w węźle połączonym z dodatnim zaciskiem źródła.

\begin{lstlisting}[language=Netlist,caption={Netlista dzielnika napięcia użytego do analizy punktu pracy},label={lst:op-divider-netlist}]
R 1 2 1k
R 2 0 1k
DCVS 1 0 5
.op
\end{lstlisting}

\begin{figure}[H]
\centering
\includegraphics[width=0.95\textwidth]{OPResultView.png}
\caption{Wynik analizy punktu pracy dla dzielnika napięcia}
\label{fig:op-result-view}
\end{figure}

\begin{figure}[H]
\centering
\includegraphics[width=0.9\textwidth]{ACVSPlusR.png}
\caption{Wynik analizy czasowej prostego układu ze źródłem sinusoidalnym i rezystorem}
\label{fig:acvs-plus-r}
\end{figure}

\begin{figure}[H]
\centering
\includegraphics[width=0.9\textwidth]{EmitterFollowerResult.png}
\caption{Wynik symulacji czasowej wtórnika emiterowego}
\label{fig:emitter-follower-result}
\end{figure}

Na wykresie wtórnika emiterowego widoczne są przebiegi napięcia bazy i emitera. Napięcie emitera podąża za napięciem bazy, ale jest przesunięte o spadek napięcia baza-emiter tranzystora. Taki wynik odpowiada oczekiwanej pracy wtórnika napięciowego: układ nie odwraca fazy sygnału i przenosi jego kształt na wyjście, ale poziom napięcia wyjściowego jest ograniczony przez właściwości złącza tranzystora. Przykład ten jest ważniejszy niż prosty układ liniowy, ponieważ jednocześnie wykorzystuje źródło sinusoidalne, kondensator, analizę czasową oraz nieliniowy model BJT.

Rozdzielenie parsera i solvera upraszcza dalszy rozwój aplikacji. Parser odpowiada za przekształcenie tekstu na dane programu, natomiast solver odpowiada za interpretację komendy analizy. W przyszłości ta sama struktura może zostać rozszerzona o kolejne polecenia, na przykład analizę małosygnałową albo warianty analizy czasowej z dodatkowymi parametrami kroku i czasu końcowego.

\section{Obsługa błędów}
\label{sec:obsluga-bledow}

Obsługa błędów została podzielona na błędy parsera i błędy stempli. Parser wykrywa problemy związane ze składnią pliku: nieznany element, złą liczbę parametrów, niepoprawny typ parametru, brak komendy analizy albo podanie więcej niż jednej komendy. Każdy taki błąd zawiera numer linii, w której wystąpił problem, co ułatwia poprawienie pliku wejściowego.

Błędy stempli dotyczą wartości, które formalnie da się odczytać z pliku, ale które nie mają sensu z punktu widzenia elementu. Przykładem jest ujemny indeks węzła albo niepoprawna wartość parametru. Parser przechwytuje takie błędy i zamienia je na błąd związany z konkretną linią netlisty. Użytkownik nie musi znać wewnętrznej struktury klas elementów, aby dowiedzieć się, gdzie w pliku znajduje się problem.

Takie rozdzielenie odpowiedzialności ułatwia utrzymanie programu. Parser nie musi znać szczegółowych warunków poprawności każdego elementu, a element nie musi wiedzieć, z której linii pliku pochodzi. Wspólnym punktem styku jest mechanizm wyjątków, który przenosi informację o niepowodzeniu do warstwy odpowiedzialnej za komunikat diagnostyczny.
